% Shared by grades 2-3 and 4-5: opening rules, the sorting problem and the region-count problem.
% Each packet defines \sortmake, the first sentence of Problem 1.
A network is dots joined by links. You may move the dots and bend the links, but each link stays joined to its own two dots and never lies across a dot. A crossing is a place where one link lies across another, and no more than two links may cross at the same place. A network is untangled when it has no crossings.

\problem \sortmake{} In the box beside it, draw it untangled or write ``stuck''.

\vspace{8pt}
\netrow{1}{n1}{5.5cm}{2.6cm}{8.4cm}{5.9cm}
\vspace{14pt}
\netrow{2}{n2}{5.5cm}{2.6cm}{8.4cm}{5.9cm}
\vspace{14pt}
\netrow{3}{n3}{5.5cm}{2.6cm}{8.4cm}{5.9cm}

\newpage
\vspace*{4pt}
\netrow{4}{n4}{5.5cm}{2.6cm}{8.4cm}{6.4cm}
\vspace{22pt}
\netrow{5}{n5}{5.5cm}{2.6cm}{8.4cm}{6.4cm}
\vspace{22pt}
\netrow{6}{n6}{5.5cm}{2.6cm}{8.4cm}{6.4cm}

\newpage
\problem A region is a part of the paper fenced in by links. The outside counts as a region too. Count the dots, links and regions in each network you untangled. Find a rule connecting the three numbers that works for every untangled network in one piece.

\vspace{14pt}
\noindent\begin{minipage}[t]{5.6cm}\vspace{0pt}
\input{gen/region_example}

\vspace{12pt}
4 dots, 5 links, 3 regions
\end{minipage}\hfill
\begin{minipage}[t]{11.6cm}\vspace{0pt}
\renewcommand{\arraystretch}{2.35}
\begin{tabular}{|>{\centering\arraybackslash}p{2.3cm}|>{\centering\arraybackslash}p{2.3cm}|>{\centering\arraybackslash}p{2.3cm}|>{\centering\arraybackslash}p{2.3cm}|}
\hline
Network & Dots & Links & Regions \\ \hline
 & & & \\ \hline
 & & & \\ \hline
 & & & \\ \hline
 & & & \\ \hline
 & & & \\ \hline
 & & & \\ \hline
 & & & \\ \hline
 & & & \\ \hline
\end{tabular}
\end{minipage}

\vspace{12pt}
\writelines{8}{\linewidth}
